% 第5章图：配置空间与局部作用域。数据：F=phi12+phi23，相邻取值不同时成本为1。
% 000:0 001:1 010:2 011:1 100:1 101:2 110:1 111:0（scratchpad ch5 验算）。
\begin{tikzpicture}[x=1mm,y=1mm,>=Stealth,
  cfg/.style={rectangle,rounded corners=1.8pt,draw=FigureStroke,line width=.7pt,fill=white,minimum width=9mm,minimum height=5mm,inner sep=1pt},
  zero/.style={cfg,fill=FigureBlueFill,draw=FigureBlue},
  two/.style={cfg,fill=FigureYellowFill,draw=FigureYellow},
  var/.style={circle,draw=FigureStroke,fill=FigureBlueFill,line width=.8pt,minimum size=6.5mm,inner sep=0pt},
  fac/.style={rectangle,draw=FigureStroke,fill=FigureGreenFill,line width=.8pt,minimum size=6mm,inner sep=1pt},
  cost/.style={font=\FigureNoteFont\fontsize{7.5}{9}\selectfont,text=FigureMuted}]
% cube: x1 -> right, x2 -> depth (up-right), x3 -> up
\def\dx{26}\def\dy{24}\def\ex{12}\def\ey{9}
\foreach \a/\b/\c/\s/\f in {0/0/0/zero/0,0/0/1/cfg/1,0/1/0/two/2,0/1/1/cfg/1,1/0/0/cfg/1,1/0/1/two/2,1/1/0/cfg/1,1/1/1/zero/0}{
  \coordinate (p\a\b\c) at ({\a*\dx+\b*\ex},{\c*\dy+\b*\ey});
}
\foreach \u/\v in {000/100,000/010,000/001,100/110,100/101,010/110,010/011,001/101,001/011,110/111,101/111,011/111}{
  \draw[FigureGrid,line width=.8pt] (p\u)--(p\v);
}
\foreach \a/\b/\c/\s/\f/\pos in {0/0/0/zero/0/below,0/0/1/cfg/1/above,0/1/0/two/2/below,0/1/1/cfg/1/above,1/0/0/cfg/1/below,1/0/1/two/2/above,1/1/0/cfg/1/below,1/1/1/zero/0/above}{
  \node[\s] (n\a\b\c) at (p\a\b\c) {$\a\b\c$};
  \ifnum\c=0 \node[cost,below=.6mm of n\a\b\c] {$F=\f$};\else \node[cost,above=.6mm of n\a\b\c] {$F=\f$};\fi
}
\draw[->,FigureMuted,line width=.6pt] (-18,-4)--(-9,-4) node[right,cost]{$x_1$};
\draw[->,FigureMuted,line width=.6pt] (-18,-4)--(-12,0.5) node[above right=-1mm,cost]{$x_2$};
\draw[->,FigureMuted,line width=.6pt] (-18,-4)--(-18,6) node[above,cost]{$x_3$};
% factor graph on the right
\begin{scope}[xshift=52mm,yshift=8mm]
 \node[var] (x1) at (0,0) {$x_1$};
 \node[var] (x2) at (16,0) {$x_2$};
 \node[var] (x3) at (32,0) {$x_3$};
 \node[fac] (f12) at (8,14) {$\phi_{1,2}$};
 \node[fac] (f23) at (24,14) {$\phi_{2,3}$};
 \draw[FigureStroke,line width=.8pt] (f12)--(x1) (f12)--(x2) (f23)--(x2) (f23)--(x3);
 \node[cost,align=center] at (16,-10) {$\phi_{1,2}$只读$x_1,x_2$\\$\phi_{2,3}$只读$x_2,x_3$};
\end{scope}
\end{tikzpicture}
